\documentclass[10pt,a4paper]{amsart}
\usepackage{amsmath,amssymb,amsthm,graphicx}
% Retained minimal preamble for the Stacks Project chapter coherent.tex (Stacks chapter, pre-delegation retained fragment for stacks-coherent).
% Verbatim excerpt of the chapter preamble (sources/stacks-spaces-descent/preamble.tex): the reference/slogan/history
% environment block (native lines 7-12), Xy-pic setup (native lines 14-19) and the native theorem declarations with
% the mathematical macros (native lines 154-217). Copied mechanically; nothing added, removed or redefined. The
% wrapper class and the amsmath/amssymb/amsthm/graphicx base packages are supplied by the assembler, so the
% class/xr-hyper/hyperref/externaldocument part of the native preamble is omitted.
% For dealing with references we use the comment environment
\usepackage{verbatim}
\newenvironment{reference}{\comment}{\endcomment}
%\newenvironment{reference}{}{}
\newenvironment{slogan}{\comment}{\endcomment}
\newenvironment{history}{\comment}{\endcomment}
% For commutative diagrams we use Xy-pic
\usepackage[all]{xy}

% We use 2cell for 2-commutative diagrams.
\xyoption{2cell}
\UseAllTwocells
% Theorem environments.
%
\theoremstyle{plain}
\newtheorem{theorem}[subsection]{Theorem}
\newtheorem{proposition}[subsection]{Proposition}
\newtheorem{lemma}[subsection]{Lemma}

\theoremstyle{definition}
\newtheorem{definition}[subsection]{Definition}
\newtheorem{example}[subsection]{Example}
\newtheorem{exercise}[subsection]{Exercise}
\newtheorem{situation}[subsection]{Situation}

\theoremstyle{remark}
\newtheorem{remark}[subsection]{Remark}
\newtheorem{remarks}[subsection]{Remarks}

\numberwithin{equation}{subsection}

% Macros
%
\def\lim{\mathop{\mathrm{lim}}\nolimits}
\def\colim{\mathop{\mathrm{colim}}\nolimits}
\def\Spec{\mathop{\mathrm{Spec}}}
\def\Hom{\mathop{\mathrm{Hom}}\nolimits}
\def\Ext{\mathop{\mathrm{Ext}}\nolimits}
\def\SheafHom{\mathop{\mathcal{H}\!\mathit{om}}\nolimits}
\def\SheafExt{\mathop{\mathcal{E}\!\mathit{xt}}\nolimits}
\def\Sch{\mathit{Sch}}
\def\Mor{\mathop{\mathrm{Mor}}\nolimits}
\def\Ob{\mathop{\mathrm{Ob}}\nolimits}
\def\Sh{\mathop{\mathit{Sh}}\nolimits}
\def\NL{\mathop{N\!L}\nolimits}
\def\CH{\mathop{\mathrm{CH}}\nolimits}
\def\proetale{{pro\text{-}\acute{e}tale}}
\def\etale{{\acute{e}tale}}
\def\QCoh{\mathit{QCoh}}
\def\Ker{\mathop{\mathrm{Ker}}}
\def\Im{\mathop{\mathrm{Im}}}
\def\Coker{\mathop{\mathrm{Coker}}}
\def\Coim{\mathop{\mathrm{Coim}}}

% Boxtimes
%
\DeclareMathSymbol{\boxtimes}{\mathbin}{AMSa}{"02}

%
% Macros for moduli stacks/spaces
%
\def\QCohstack{\mathcal{QC}\!\mathit{oh}}
\def\Cohstack{\mathcal{C}\!\mathit{oh}}
\def\Spacesstack{\mathcal{S}\!\mathit{paces}}
\def\Quotfunctor{\mathrm{Quot}}
\def\Hilbfunctor{\mathrm{Hilb}}
\def\Curvesstack{\mathcal{C}\!\mathit{urves}}
\def\Polarizedstack{\mathcal{P}\!\mathit{olarized}}
\def\Complexesstack{\mathcal{C}\!\mathit{omplexes}}
% \Pic is the operator that assigns to X its picard group, usage \Pic(X)
% \Picardstack_{X/B} denotes the Picard stack of X over B
% \Picardfunctor_{X/B} denotes the Picard functor of X over B
\def\Pic{\mathop{\mathrm{Pic}}\nolimits}
\def\Picardstack{\mathcal{P}\!\mathit{ic}}
\def\Picardfunctor{\mathrm{Pic}}
\def\Deformationcategory{\mathcal{D}\!\mathit{ef}}
\usepackage{cleveref}
\crefname{originallemma}{original lemma}{original lemmas}
\Crefname{originallemma}{Original lemma}{Original lemmas}
\makeatletter
\newlabel{algebra-lemma-S-plus-generated}{{07Z4}{source}{}{sourcepointer-1}{}}
\newlabel{algebra-lemma-S-plus-generated@cref}{{[originallemma][0][]07Z4}{[1][0][]source}{}{}{}}
\makeatother
\crefname{originallemma}{original lemma}{original lemmas}
\Crefname{originallemma}{Original lemma}{Original lemmas}
\makeatletter
\newlabel{algebra-lemma-dehomogenize-finite-type}{{07Z2}{source}{}{sourcepointer-2}{}}
\newlabel{algebra-lemma-dehomogenize-finite-type@cref}{{[originallemma][0][]07Z2}{[1][0][]source}{}{}{}}
\makeatother
\crefname{originallemma}{original lemma}{original lemmas}
\Crefname{originallemma}{Original lemma}{Original lemmas}
\makeatletter
\newlabel{algebra-lemma-graded-Noetherian}{{00JW}{source}{}{sourcepointer-3}{}}
\newlabel{algebra-lemma-graded-Noetherian@cref}{{[originallemma][0][]00JW}{[1][0][]source}{}{}{}}
\makeatother
\crefname{originallemma}{original lemma}{original lemmas}
\Crefname{originallemma}{Original lemma}{Original lemmas}
\makeatletter
\newlabel{constructions-equation-global-sections-more-generally}{{01MR}{source}{}{sourcepointer-4}{}}
\newlabel{constructions-equation-global-sections-more-generally@cref}{{[originallemma][0][]01MR}{[1][0][]source}{}{}{}}
\makeatother
\crefname{originallemma}{original lemma}{original lemmas}
\Crefname{originallemma}{Original lemma}{Original lemmas}
\makeatletter
\newlabel{constructions-equation-multiply-more-generally}{{01MQ}{source}{}{sourcepointer-5}{}}
\newlabel{constructions-equation-multiply-more-generally@cref}{{[originallemma][0][]01MQ}{[1][0][]source}{}{}{}}
\makeatother
\crefname{originallemma}{original lemma}{original lemmas}
\Crefname{originallemma}{Original lemma}{Original lemmas}
\makeatletter
\newlabel{constructions-lemma-proj-sheaves}{{01M7}{source}{}{sourcepointer-6}{}}
\newlabel{constructions-lemma-proj-sheaves@cref}{{[originallemma][0][]01M7}{[1][0][]source}{}{}{}}
\makeatother
\crefname{originallemma}{original lemma}{original lemmas}
\Crefname{originallemma}{Original lemma}{Original lemmas}
\makeatletter
\newlabel{constructions-lemma-where-invertible}{{01MU}{source}{}{sourcepointer-7}{}}
\newlabel{constructions-lemma-where-invertible@cref}{{[originallemma][0][]01MU}{[1][0][]source}{}{}{}}
\makeatother
\crefname{originallemma}{original lemma}{original lemmas}
\Crefname{originallemma}{Original lemma}{Original lemmas}
\makeatletter
\newlabel{lemma-coherent-on-proj-general}{{0B5Q}{source}{}{sourcepointer-8}{}}
\newlabel{lemma-coherent-on-proj-general@cref}{{[originallemma][0][]0B5Q}{[1][0][]source}{}{}{}}
\makeatother
\crefname{originallemma}{original lemma}{original lemmas}
\Crefname{originallemma}{Original lemma}{Original lemmas}
\makeatletter
\newlabel{properties-lemma-proj-quasi-coherent}{{0AG5}{source}{}{sourcepointer-9}{}}
\newlabel{properties-lemma-proj-quasi-coherent@cref}{{[originallemma][0][]0AG5}{[1][0][]source}{}{}{}}
\makeatother
\title{Recovering the tail of a graded module over a general Noetherian graded ring: for A Noetherian graded, d the lcm of degrees of generators over A\_0, M a finite graded module, X = Proj(A), k in Z: N' = bigoplus\_\{n >= k\} H\textasciicircum{}0(X, M(n)\textasciicircum{}tilde) and N = bigoplus\_\{n >= k\} H\textasciicircum{}0(X, M\textasciicircum{}tilde(n)) are finite A-modules, there are canonical maps N -> N' and (k small) M -> N', M\_n -> N'\_n is an isomorphism for n >> 0, and N\_n -> N'\_n is an isomorphism whenever d divides n.}
\author{The Stacks Project Authors}
\date{stacks-project:coherent.tex:local-snapshot-20261012; GNU FDL 1.2 or later}
\begin{document}
\maketitle
\noindent\textit{Editorial note.} Statement span is lines 4297-4314 of sources/stacks-coherent/pinned-original-source.tex; proof span is lines 4316-4414, the complete author proof (97 body lines).\par\medskip
\noindent\textit{Editorial note.} Navigation stubs: 9 cited-but-not-included labels are supplied as non-load-bearing pointers with their pinned tags where the snapshot tag index has them.\par\medskip
\noindent\textit{Editorial note.} extra\_packages is empty on purpose: the retained preamble already loads xy with options and a second declaration clashes in the assembler wrapper.\par\medskip
\noindent\textit{Editorial note.} Text is retained verbatim, including the author's own typos and placeholders.\par\medskip
\noindent\textit{Editorial note.} Card-surface disclosure per the tightened ruling: the identity of the composite N\textasciicircum{}tilde -> M\textasciicircum{}tilde -> N'\textasciicircum{}tilde with the constructed isomorphism is the author's named omission ('proof omitted; hint'), and shift-compatibility of N' is called immediate from the construction; all six enumerated claims carry their proofs in the span.\par\medskip
\noindent\textit{Editorial note.} Span computed from the \textbackslash{}label marker through the enclosing \textbackslash{}end\{lemma\}; precard\_gate RESULT=OK (probe L9999).\par\medskip
\noindent\textit{Allowed standard prerequisites (editorial).}\begin{itemize}
\item Explicit prerequisite boundary. See why-field; every pointer named here occurs verbatim in the retained statement/proof text.
\end{itemize}\medskip
\ifdefined\hypertarget\hypertarget{sourcepointer-1}{}\fi
\noindent\textit{Source reference.} Original lemma 07Z4: sources/stacks-coherent/pinned-original-source.tex, source lines 4414--4414. Cited by the retained statement or proof but not part of the collected spans; supplied only as a navigation pointer. The argument does not depend on its content.\par\medskip
\ifdefined\hypertarget\hypertarget{sourcepointer-2}{}\fi
\noindent\textit{Source reference.} Original lemma 07Z2: sources/stacks-coherent/pinned-original-source.tex, source lines 4414--4414. Cited by the retained statement or proof but not part of the collected spans; supplied only as a navigation pointer. The argument does not depend on its content.\par\medskip
\ifdefined\hypertarget\hypertarget{sourcepointer-3}{}\fi
\noindent\textit{Source reference.} Original lemma 00JW: sources/stacks-coherent/pinned-original-source.tex, source lines 4414--4414. Cited by the retained statement or proof but not part of the collected spans; supplied only as a navigation pointer. The argument does not depend on its content.\par\medskip
\ifdefined\hypertarget\hypertarget{sourcepointer-4}{}\fi
\noindent\textit{Source reference.} Original lemma 01MR: sources/stacks-coherent/pinned-original-source.tex, source lines 4414--4414. Cited by the retained statement or proof but not part of the collected spans; supplied only as a navigation pointer. The argument does not depend on its content.\par\medskip
\ifdefined\hypertarget\hypertarget{sourcepointer-5}{}\fi
\noindent\textit{Source reference.} Original lemma 01MQ: sources/stacks-coherent/pinned-original-source.tex, source lines 4414--4414. Cited by the retained statement or proof but not part of the collected spans; supplied only as a navigation pointer. The argument does not depend on its content.\par\medskip
\ifdefined\hypertarget\hypertarget{sourcepointer-6}{}\fi
\noindent\textit{Source reference.} Original lemma 01M7: sources/stacks-coherent/pinned-original-source.tex, source lines 4414--4414. Cited by the retained statement or proof but not part of the collected spans; supplied only as a navigation pointer. The argument does not depend on its content.\par\medskip
\ifdefined\hypertarget\hypertarget{sourcepointer-7}{}\fi
\noindent\textit{Source reference.} Original lemma 01MU: sources/stacks-coherent/pinned-original-source.tex, source lines 4414--4414. Cited by the retained statement or proof but not part of the collected spans; supplied only as a navigation pointer. The argument does not depend on its content.\par\medskip
\ifdefined\hypertarget\hypertarget{sourcepointer-8}{}\fi
\noindent\textit{Source reference.} Original lemma 0B5Q: sources/stacks-coherent/pinned-original-source.tex, source lines 4237--4255. Cited by the retained statement or proof but not part of the collected spans; supplied only as a navigation pointer. The argument does not depend on its content.\par\medskip
\ifdefined\hypertarget\hypertarget{sourcepointer-9}{}\fi
\noindent\textit{Source reference.} Original lemma 0AG5: sources/stacks-coherent/pinned-original-source.tex, source lines 4414--4414. Cited by the retained statement or proof but not part of the collected spans; supplied only as a navigation pointer. The argument does not depend on its content.\par\medskip
\begin{lemma}
\label{lemma-recover-tail-graded-module-general}
Let $A$ be a Noetherian graded ring and let $d$ be the lcm of generators
of $A$ over $A_0$. Let $M$ be a finite graded $A$-module.
Set $X = \text{Proj}(A)$ and let $\widetilde{M}$ be
the quasi-coherent $\mathcal{O}_X$-module on $X$ associated to $M$.
Let $k \in \mathbf{Z}$.
\begin{enumerate}
\item $N' = \bigoplus_{n \geq k} H^0(X, \widetilde{M(n)})$
is a finite $A$-module,
\item $N = \bigoplus_{n \geq k} H^0(X, \widetilde{M}(n))$
is a finite $A$-module,
\item there is a canonical map $N \to N'$,
\item if $k$ is small enough there is a canonical map $M \to N'$,
\item the map $M_n \to N'_n$ is an isomorphism for $n \gg 0$,
\item $N_n \to N'_n$ is an isomorphism for $d | n$.
\end{enumerate}
\end{lemma}

\begin{proof}
The map $N \to N'$ in (3) comes from
Constructions, Equation (\ref{constructions-equation-multiply-more-generally})
by taking global sections.

\medskip\noindent
By
Constructions, Equation
(\ref{constructions-equation-global-sections-more-generally})
there is a map of graded $A$-modules
$M \to \bigoplus_{n \in \mathbf{Z}} H^0(X, \widetilde{M(n)})$.
If the generators of $M$ sit in degrees $\geq k$, then the image
is contained in the submodule
$N' \subset \bigoplus_{n \in \mathbf{Z}} H^0(X, \widetilde{M(n)})$
and we get the map in (4).

\medskip\noindent
By Algebra, Lemmas \ref{algebra-lemma-graded-Noetherian} and
\ref{algebra-lemma-S-plus-generated} the ring $A_0$ is Noetherian
and $A$ is generated over $A_0$ by finitely many elements
$f_1, \ldots, f_r$ homogeneous of positive degree.
Let $d = \text{lcm}(\deg(f_i))$. Then we see that (6) holds
for example by
Constructions, Lemma \ref{constructions-lemma-where-invertible}.

\medskip\noindent
Because $M$ is a finite $A$-module we see that
$\widetilde{M}$ is a finite type $\mathcal{O}_X$-module,
i.e., a coherent $\mathcal{O}_X$-module.
Thus part (2) follows from Lemma \ref{lemma-coherent-on-proj-general}.

\medskip\noindent
We will deduce (1) from (2) using a trick. For $q \in \{0, \ldots, d - 1\}$
write
$$
{}^qN = \bigoplus\nolimits_{n + q \geq k} H^0(X, \widetilde{M(q)}(n))
$$
By part (2) these are finite $A$-modules. The Noetherian ring $A$
is finite over $A^{(d)} = \bigoplus_{n \geq 0} A_{dn}$, because
it is generated by $f_i$ over $A^{(d)}$ and $f_i^d \in A^{(d)}$.
Hence ${}^qN$ is a finite $A^{(d)}$-module.
Moreover, $A^{(d)}$ is Noetherian (follows from
Algebra, Lemma \ref{algebra-lemma-dehomogenize-finite-type}).
It follows that the $A^{(d)}$-submodule
${}^qN^{(d)} = \bigoplus_{n \in \mathbf{Z}} {}^qN_{dn}$
is a finite module over $A^{(d)}$. Using the isomorphisms
$\widetilde{M(dn + q)} = \widetilde{M(q)}(dn)$ we can write
$$
N' =
\bigoplus\nolimits_{q \in \{0, \ldots, d - 1\}}
\bigoplus\nolimits_{dn + q \geq k}
H^0(X, \widetilde{M(q)}(dn)) =
\bigoplus\nolimits_{q \in \{0, \ldots, d - 1\}} {}^qN^{(d)}
$$
Thus $N'$ is finite over $A^{(d)}$ and a fortiori finite over $A$.
Thus (1) is true.

\medskip\noindent
Let $K$ be a finite $A$-module such that $\widetilde{K} = 0$. We claim
that $K_n = 0$ for $d|n$ and $n \gg 0$. Arguing as above we see that
$K^{(d)}$ is a finite $A^{(d)}$-module. Let
$x_1, \ldots, x_m \in K$ be homogeneous generators of $K^{(d)}$
over $A^{(d)}$, say sitting in degrees $d_1, \ldots, d_m$ with $d | d_j$.
For each $i$ and $j$ there exists an $n_{ij} \geq 0$ such that
$f_i^{n_{ij}} x_j = 0$ in $K_{d_j + n_{ij}}$: if not then
$x_j/f_i^{d_i/\deg(f_i)} \in K_{(f_i)}$ would not be zero, i.e.,
$\widetilde{K}$ would not be zero. Here we use that $\deg(f_i) | d | d_j$
for all $i, j$. We conclude that $K_n$ is zero for
$n$ with $d | n$ and $n > \max_j (d_j + \sum_i n_{ij} \deg(f_i))$
as every element of $K_n$ is a sum of terms where each term is a
monomials in the $f_i$ times one of the $x_j$ of total degree $n$.

\medskip\noindent
To finish the proof, we have to show that $M \to N'$ is an isomorphism
in all sufficiently large degrees. The map $N \to N'$ induces an
isomorphism $\widetilde{N} \to \widetilde{N'}$ because on the affine
opens $D_+(f_i) = D_+(f_i^d)$ the corresponding modules are isomorphic:
$N_{(f_i)} \cong N_{(f_i^d)} \cong N'_{(f_i^d)} \cong N'_{(f_i)}$
by property (6).
By Properties, Lemma \ref{properties-lemma-proj-quasi-coherent}
we have a canonical isomorphism $\widetilde{N} \to \widetilde{M}$.
The composition $\widetilde{N} \to \widetilde{M} \to \widetilde{N'}$
is the isomorphism above (proof omitted; hint: look on standard
affine opens to check this). Thus the map $M \to N'$ induces an isomorphism
$\widetilde{M} \to \widetilde{N'}$.
Let $K = \Ker(M \to N')$ and $Q = \Coker(M \to N')$.
Recall that the functor
$M \mapsto \widetilde{M}$ is exact, see
Constructions, Lemma \ref{constructions-lemma-proj-sheaves}.
Hence we see that $\widetilde{K} = 0$ and $\widetilde{Q} = 0$.
By the result of the previous paragraph we see that $K_n = 0$ and
$Q_n = 0$ for $d | n$ and $n \gg 0$. At this point we finally see
the advantage of using $N'$ over $N$: the functor $M \leadsto N'$
is compatible with shifts (immediate from the construction).
Thus, repeating the whole argument with $M$ replaced by $M(q)$
we find that $K_n = 0$ and $Q_n = 0$ for $n \equiv q \bmod d$
and $n \gg 0$. Since there are
only finitely many congruence classes modulo $n$ the proof is finished.
\end{proof}

\end{document}
